Class imbalance is usually handled by resampling the training data and then fitting an ensemble. Most hybrid
pipelines apply one resampler to the whole minority class and combine their learners with fixed weights. We
propose \ars{} (Adaptive Resampler-Selection Stacking), which clusters the minority class, scores each cluster
for class overlap and local density, selects a resampler \emph{per cluster}, and combines Random Forest, XGBoost and
LightGBM with confidence weights. We evaluate it under a leakage-free, budget-matched protocol on 28 datasets
(credit-card fraud, mammography, 10-class yeast and 25 benchmarks) against eight baselines and ablations: 3,780
cross-validation folds, analysed with frequentist and Bayesian tests. \ars{} is significantly and practically better
than global SMOTE-ENN (Bayesian probability $\geq 0.96$ on F1, AUC-PR and MCC). It is, however, \emph{practically
equivalent} to plain SMOTE, SMOTE-Tomek and its own ablations, and it cannot be separated from training without
resampling. Diagnostics explain why: the confidence weights are almost uniform, Tomek cleaning removes nothing after
1:1 SMOTE, and the selector rarely varies its choice within a dataset. A pilot of validated selection closes the gap
to the baselines but does not exceed them. Resampling helps mainly when the minority class is very small (mean
AUC-PR gain $+0.012$ below 75 minority samples, $-0.009$ above). We release the reproducible pipeline and discuss
five evaluation pitfalls that make new resamplers look better than they are.
